\noindent
The main objectives of the project are:

\begin{itemize}
    \item Develop high-quality 3D digital representations of cultural heritage sites
    \item Enable virtual access through VR/AR technologies
    \item Reduce physical tourism pressure on vulnerable sites
    \item Establish standardized methodologies for digital heritage documentation
    \item Foster interdisciplinary collaboration between technology and humanities
    \item Support open science through accessible datasets and publications
\end{itemize}

\vspace{1em}
\noindent In detail, the project aims to save lesser-known or rapidly deteriorating monuments using 3D scanning and VR/AR technology, allowing them to be visited without the physical burden of mass tourism. Key targeted achievements include:
\begin{itemize}
    \item \textbf{Technological Validation:} Testing XR platforms in complex, outdoor environments (e.g., ruins, built heritage) to increase Technology Readiness Levels (TRL) and demonstrate real-world impact.
    \item \textbf{Open Science and Visibility:} Generating high-quality datasets, 3D models, and VR assets as part of an open access ecosystem, maximizing visibility, citation, and future funding opportunities.
    \item \textbf{International Network Expansion:} Strengthening EU consortium networks by bringing together technology hubs, scientific publishers, and top-tier heritage protection partners across the Mediterranean, Central-Eastern, and Atlantic European regions.
\end{itemize}